\documentclass[10pt,a4paper]{article}
\usepackage[margin=22mm]{geometry}
\usepackage[T1]{fontenc}
\usepackage{lmodern,amsmath,amssymb,booktabs,graphicx,tikz,hyperref,fvextra}
\usetikzlibrary{arrows.meta,positioning,calc,decorations.pathmorphing}
\definecolor{ink}{RGB}{20,53,78}
\hypersetup{colorlinks=true,urlcolor=ink,linkcolor=ink}
\setlength{\parindent}{0pt}
\setlength{\parskip}{5pt}
\emergencystretch=3em

\title{\color{ink}Control of Dynamical Systems: An Original Learning Path}
\author{}
\date{Original learning collection -- November 2026}
\begin{document}
\maketitle
\textbf{Evidence note.} Original teaching scenario. Numerical parameters are deliberately chosen for learning. All generated data and curves are synthetic; no physical measurement is asserted.
\section{Purpose and prerequisites}
This original learning path connects physical balances to control decisions. It assumes elementary differential equations, Laplace transforms, complex numbers, linear algebra and basic Python. Its examples are deliberately hypothetical so that parameter choices, model limits and reproducibility remain visible.
\section{Learning outcomes}
After completing the collection, a learner should be able to identify states and inputs; distinguish physical states from measurements; linearize a nonlinear balance; derive and audit transfer functions; test stability algebraically and numerically; design classical and state-space feedback; implement a sampled controller; and separate model fitting from validation.
\section{A four-stage sequence}
\begin{enumerate}
\item Build and inspect models: thermal enclosure, ventilation, vibration isolation and passive notch filtering.
\item Close feedback loops: motor load rejection, root-locus gain limits and communication-delay budgets.
\item Design with explicit tradeoffs: lead-lag accuracy, actuator-limited PID and observer-based thermal regulation.
\item Check implementation and evidence: sampling stability and synthetic process identification with held-out excitation.
\end{enumerate}
\section{A suggested study rhythm}
For each laboratory, first solve the balance or characteristic equation without running code. Predict one qualitative result and one numerical limit. Then execute the notebook, compare with the prediction, and write a short account of assumptions and discrepancies. Each laboratory has five integrated problems with three linked parts. Begin with the physical illustration where provided, then derive the model, calculate its behavior and make a justified design decision. The general collection adds five cross-topic cases with six parts each, including interacting water tanks. The separate worked solutions expose intermediate steps and infeasible requirements.
\newpage
\section{First pass and return tasks}
The numbered laboratories form a coherent topic sequence, not a promise that every exercise is harder than every earlier one. In Laboratory 02, first solve Problems 1--3 and 5(b), then return to feedback Problems 4 and 5(a,c) after Laboratory 05. In Laboratory 05, revisit the Routh items 4(a,c) after Laboratory 06 and controllability item 1(b) after Laboratory 10. The sensor-noise calculation in Laboratory 01 and the statistical interval in Laboratory 12 assume elementary probability.
For the general collection, the recommended first-pass order is Case 2 (vibration), Case 1 (tanks), Case 3 (motor), Case 4 (camera), then Case 5 (thermal estimation and sampling). Begin Case 2 after Laboratory 03; attempt Cases 1 and 3 after feedback and stability; complete Case 4 after Laboratory 08 and Case 5 after Laboratory 11. Return to validation-plan parts after Laboratory 12. The repository's learning-path page gives explicit prerequisites and exit checks.

\section{A self-assessment rubric}
Assess each submission on four dimensions: correct physical interpretation and units; a traceable derivation; reproducible computations with purposeful checks; and conclusions proportional to the evidence. A visually smooth simulation with an incorrect input location is not a correct model. A passing assertion is not a physical experiment.
\section{Suggested extensions}
Introduce uncertainty one parameter at a time; compare analytical formulas with direct integration; change a disturbance entry point; add actuator or sensor dynamics; and test whether a controller remains stable after discretization. Collect real data only in a separately documented and appropriately supervised experimental setting.
\section{Authorship boundary}
The teaching scenarios, assignment wording, derivations, figures and programs were newly prepared for this repository. Supplied course materials served only to identify broad subjects. No source worksheets, translated question lists, copied figures or original administrative material are distributed. Established scientific methods are cited and are not claimed as novel inventions.
\section{Method references}
The equations use standard balance laws and feedback-control methods. Further explanation: \url{https://apmonitor.com/pdc/} and \url{https://fbsbook.org/}. Implementation documentation: \url{https://python-control.readthedocs.io/en/0.10.2/}. See the repository reference catalog for individual topics and videos.
\end{document}
